\documentclass{article}
\usepackage{amsmath}
\usepackage[hidelinks]{hyperref}
\title{One-byte terminal evidence: timing, provisioning, and old mint authority}
\date{}
\begin{document}
\maketitle

\paragraph{Design and comparison boundary.}
The civil source, keyed simulator publication and native radio
accounting are the selected and pinned fragments used in the joined
closure replay. This is a held-in case using the same task cut, not
a new mission. No real protocol packet, radio-header, cryptographic
authentication or physical network cost is inferred from payload bytes.
The trusted callbacks, two configuration messages, decoder state,
old application queue and mint authority are explicit adapter laws.
Use \verb|SOURCE.tsv| and the two patch files to recover the test.
The only receiver-side state inspection is of its local binding.
The canonical station applies the same immutable job before emitting
a matching result and cannot know when the receiver silently loses
its context.

\paragraph{Completed and pending look different only through valid evidence.}
When the published old result is durable and the source still has
the exact prebound dictionary, the single code \texttt{s} names that
result. In the masked case Base has four bytes already charged and
can send one, not the three-character full receipt. At tick six a
real one-byte callback lets the original carrier physically reenter
and the optional transfer safely use Q's backup capacity. Sending
full precision instead triggers two more station attempts and a
later three-byte terminal return. The same-key late control retry
still runs in both configurations, and both independent side jobs run.

Across both actual station outcomes, loss of the binding deletes the
decoder row: old duty continues via fallback, not via an assumed
agreement. In the reassignment mode an older receipt has legitimately
been minted for a committed predecessor, queued, and sent after the
local row has been rewritten to the new duty. The old and new origin
can yield the very same code, source, endpoint, current mapping and
time. The equally visible still-pending run cannot safely release.
The sound decoder regards both origins as still possible and declines
both short receipts until authoritative full evidence arrives; nothing
about the callback certifies which mapping was in force at mint time.
Treating the current slot map as the old one releases early and
strands the still-pending results until after the deadline.

\paragraph{Exact retirement rule, not a TTL heuristic.}
For a visible code at a cut, the lawful terminal inference is
universal over admissible mint and travel histories compatible with
that code, authority, endpoint and receiver state. All histories must
entail the claimed settlement. A unique job key is sufficient, but
distinct origins that entail the very same terminal fact need not
remain distinct for this check. Let the last possible old mint
be $m$ and let $L$ bound all remaining time in flight, including
queues and retransmission of already minted bytes. Replay minting
abilities that aren't retired cannot be treated as already in flight.
On one discrete clock and with both interval endpoints attainable
jointly, an old receipt can still land at or after reopening $r$
exactly when $r\leq m+L$. The latest mint with longest delay witnesses
the equality case, so reuse requires $r>m+L$. If only a joint
nonrectangular envelope is given, compute the reachable arrival
support instead of declaring its separated maxima attainable.
If neither old mint retirement nor a flight bound can be justified,
this argument cannot authorize a finite timeout. Use a disjoint
receipt context or a full self-identifying result. The interval
calculation is independently checked in \verb|Epoch_Fence.lean|;
it does not establish those premises for the stock radio or real plants.

\paragraph{Cost and null models.}
The masked committed witness goes from 23 to 18 debited operational
payload bytes, and from tick-eight to tick-six source release. Those
totals include the optional callback in the compact arm, the
retained late duplicate, and separate side traffic. The masked
unpublished partner stays at 23 operating bytes. Context installation
is an actual four-byte delivered configuration and a two-byte
acknowledgement before the test window. Its six bytes make the
per-use total in the first witness 24 rather than 18; advance
provisioning is not free. Restoring a second mapping in the
reassignment arm charges a second complete installation. Both the
binding-lost and reassigned sound arms use the full late receipt
and still meet the old deadline in all test histories. Relative
to the full arm they do not save traffic.

The eight-replay sweep for each arm varies worker outcome, common
interruption and independent primary path failure; these are not
weighted estimates. The tight matched comparator is a provisioned
static dictionary plus keyed retry. It ties trace for trace. The
design abstracts from (but does not implement)
\href{https://www.rfc-editor.org/rfc/rfc8724.html#section-6}{SCHC
scoped rule identifiers and endpoint context} and
\href{https://www.rfc-editor.org/rfc/rfc9175.html#section-4}{CoAP
request-response binding}; compare
\href{https://web.stanford.edu/~ouster/cgi-bin/papers/rifl.pdf}{RIFL
Section 3} for retaining completion records and retry rendezvous.
These precedents make the full-content arm a timing and cost
calibration, not a strongest-prior advance. Amortization requires
a separate workload that earns stable cross-operation context; per-use
installation at matched scope loses the measured byte saving here.
No change to independent transfer, strongest-baseline or coverage status.
\end{document}
